\documentclass[10pt]{amsart}

\usepackage[T1]{fontenc}
\usepackage{lmodern}
\usepackage[protrusion=true,expansion=false]{microtype}
\usepackage{mathtools}
\usepackage{amssymb}
\usepackage{booktabs}
\usepackage{enumitem}
\usepackage{xcolor}
\usepackage[colorlinks=true,linkcolor=blue!45!black,citecolor=blue!45!black,urlcolor=blue!55!black]{hyperref}
\usepackage[nameinlink,capitalize,noabbrev]{cleveref}

\numberwithin{equation}{section}

\newtheorem{theorem}{Theorem}[section]
\newtheorem{proposition}[theorem]{Proposition}
\newtheorem{lemma}[theorem]{Lemma}
\newtheorem{corollary}[theorem]{Corollary}
\theoremstyle{definition}
\newtheorem{definition}[theorem]{Definition}
\theoremstyle{remark}
\newtheorem{remark}[theorem]{Remark}

\newcommand{\CC}{\mathbf C}
\newcommand{\QQ}{\mathbf Q}
\newcommand{\PP}{\mathbf P}
\newcommand{\GGm}{\mathbf G_m}
\newcommand{\IC}{\operatorname{IC}}
\newcommand{\Nm}{\operatorname{Nm}}
\newcommand{\Alt}{\operatorname{Alt}}
\newcommand{\Sym}{\operatorname{Sym}}
\newcommand{\Stab}{\operatorname{Stab}}
\newcommand{\Sing}{\operatorname{Sing}}
\newcommand{\id}{\operatorname{id}}
\newcommand{\reg}{\mathrm{reg}}
\newcommand{\exc}{\mathrm{exc}}
\newcommand{\cH}{\mathcal H}
\newcommand{\cI}{\mathcal I}
\newcommand{\cM}{\mathcal M}
\newcommand{\cN}{\mathcal N}
\newcommand{\cB}{\mathcal B}
\newcommand{\cP}{\mathcal P}
\newcommand{\cX}{\mathcal X}
\newcommand{\cR}{\mathcal R}
\newcommand{\cU}{\mathcal U}
\newcommand{\cA}{\mathcal A}
\newcommand{\cC}{\mathcal C}
\newcommand{\cL}{\mathcal L}
\newcommand{\wh}{\widehat}
\newcommand{\wt}{\widetilde}
\newcommand{\bo}{\boxtimes}

\title[A realization of $G_2$ for Litt Problem No. 13]
      {A Realization of $G_2$ for Litt Problem No. 13}
\author{OpenAI}
\date{July 26, 2026}

\hypersetup{
  pdftitle={A Realization of G2 for Litt Problem No. 13},
  pdfauthor={OpenAI}
}

\subjclass[2020]{14K12, 14H40, 18M25, 20G41}
\keywords{abelian variety, cyclic cover, Prym variety, intersection complex,
convolution, Tannaka group, exceptional group $G_2$}

\begin{document}
\raggedbottom

\begin{abstract}
We construct a complex abelian threefold $P$ and an integral symmetric ample
divisor $X\subset P$ whose full convolution Tannaka group is the simple group
of type $G_2$.  The construction starts from a cyclic triple cover
$C\to E$ of an elliptic curve, equivalently from a smooth curve
$C:w^6=f_3(x)$.  The Prym polarization has type $(1,1,3)$, and the theta
section $X$ has a unique simple elliptic singularity of type
$\wt E_7$.  Its intersection complex $K=\IC_X$ satisfies
\[
  \chi(K)=7,
  \qquad
  K\ \text{is a direct summand of }\Alt^2(K).
\]
Here the second assertion is in the convolution Tannakian category; it is
obtained from an exact universal norm--conic identity and a support calculation
for pair convolution.  The tautological representation is faithful,
orthogonal, and irreducible on the identity component; its Hodge cocharacter
has multiplicities $(2,3,2)$.  These data identify the full group, including
its component group, with $G_2$ in its standard seven-dimensional
representation.
\end{abstract}

\maketitle

\section{Introduction}

Let $A$ be a complex abelian variety.  Convolution of constructible
complexes on $A$ induces a neutral Tannakian structure after quotienting by
the negligible objects.  Thus an integral subvariety $Z\subset A$ with
nonzero intersection-cohomology Euler characteristic determines an
algebraic group $G(\IC_Z)$ and a faithful representation of dimension
$\chi(\IC_Z)$; see \cite{KraemerWeissauer2015}.  Litt Problem No.~13 asks
whether the exceptional simple groups
\[
  G_2,\qquad E_7,\qquad E_8,\qquad F_4
\]
occur in this way \cite{Litt13}.

This paper settles the $G_2$ case.  Let
\[
 \cB=\operatorname{Spec}
 \CC[A,B,(4A^3+27B^2)^{-1}].
\]
For $b=(A,B)\in\cB(\CC)$, consider the smooth projective curves
\[
 E_b:y^2=x^3+Ax+B,
 \qquad
 C_b:w^6=x^3+Ax+B,
\]
and the cyclic triple covering
\[
 \pi_b:C_b\longrightarrow E_b,
 \qquad (x,w)\longmapsto(x,w^3).
\]
Write $O\in E_b$ for the point at infinity, put
\[
 \kappa_b=\pi_b^*\mathcal O_{E_b}(O),
 \qquad
 P_b=\ker(\Nm_{\pi_b})^0,
\]
and define the cyclic theta section
\[
 X_b=\{M\in P_b:h^0(C_b,\kappa_b\otimes M)>0\}.
\]

\begin{theorem}[Main theorem]\label{thm:main}
There is a nonempty Zariski-open subset $\cB^\circ\subset\cB$ such that,
for every $b\in\cB^\circ(\CC)$, the variety $P_b$ is an abelian threefold,
$X_b\subset P_b$ is an integral symmetric ample divisor, and, as an
algebraic group together with its tautological module,
\[
  \bigl(G(\IC_{X_b}),\omega(\IC_{X_b})\bigr)\simeq(G_2,V_7),
\]
where $V_7$ is the standard faithful seven-dimensional representation.  This
is the full
convolution Tannaka group, not only its identity component or its connected
derived group.
\end{theorem}

The proof has three parts.  First, the geometry of the norm fibre gives a
minimal resolution $\rho:S\to X$ with
\[
 R\rho_*\QQ_S[2]\simeq\IC_X\oplus\delta_0,
 \qquad
 \chi(\IC_X)=7,
\]
and the generic-twist Hodge vector of $\IC_X$ is $(2,3,2)$.  Second, a
universal norm calculation shows that the generic pair-incidence curve
\[
 C_x=\{y\in X:-x-y\in X\}
\]
has precisely two nonexceptional components, globally labelled over a dense
open subset of $X$, and transposition exchanges them.  A top-cohomology
rank calculation and the decomposition theorem then give
\[
 \IC_X\subset\Alt^2(\IC_X).
\]
Finally, the Hodge cocharacter and the exterior-square inclusion distinguish
$G_2$ from the other connected irreducible orthogonal subgroups of
$O_7$.  The same exterior-square inclusion excludes the only possible
disconnected extension.

The universal norm identity is an exact calculation over
$\QQ(\lambda,e_2,e_3,u)$, not a numerical experiment.  The supplementary
verification performs rational-function arithmetic and independently checks
that a cleared polynomial numerator is identically zero.  Finite-field
calculations are used only to certify that specified Zariski-open conditions
are nonempty.

All varieties are over $\CC$.  Constructible complexes and Hodge modules
retain their canonical $\QQ$-structures.  Every Tannaka group and
tautological module below means the corresponding object after extension of
scalars from $\QQ$ to $\CC$; in particular, all algebraic groups are over
$\CC$.

\section{Tannakian conventions}

Let $A$ be a complex abelian variety and let $m:A\times A\to A$ be addition.
For constructible complexes we write
\[
 K*L=Rm_*(K\bo L).
\]
Perverse sheaves of Euler characteristic zero generate the negligible
tensor ideal.  The quotient of the semisimple perverse category by this
ideal is rigid symmetric monoidal, and the tensor subcategory generated by
one nonnegligible object is neutral Tannakian
\cite{KraemerWeissauer2015,Schnell2015}.

For an integral subvariety $Z\subset A$ of dimension $d$, let
\[
 \IC_Z=j_{!*}\QQ_{Z_{\reg}}[d]
\]
be its perverse intersection complex, extended by zero to $A$.  If
$\chi(\IC_Z)>0$, we denote by $G(\IC_Z)$ the complexification of the
Tannaka group of the category generated by $\IC_Z$, and by $\omega(\IC_Z)$
its complex tautological module.  Then
\[
 \dim\omega(\IC_Z)=\chi(\IC_Z).
\]
The action on the tautological module is faithful by construction.

The commutativity constraint contains the usual cohomological sign.  Since
the support in this paper has even dimension two, the alternating projector
on $\IC_X*\IC_X$ is carried by the fibre functor to the ordinary exterior
square of $\omega(\IC_X)$.  We write $\Alt^2(\IC_X)$ for its image in the
Tannakian quotient.

\section{The cyclic theta section}

We suppress the parameter $b$ until \cref{sec:open-family}.  Thus
\[
 E:y^2=f_3(x),\qquad C:w^6=f_3(x),\qquad y=w^3,
\]
where $f_3$ is a separable cubic, and $\pi:C\to E$ is the resulting cyclic
triple cover.

\subsection{The cover, the Prym, and the polarization}

The three branch points of $\pi$ are the nonzero points of $E[2]$.  If
$\eta=\mathcal O_E(O)$, then the cyclic covering algebra is
\[
 \pi_*\mathcal O_C
 =\mathcal O_E\oplus\eta^{-1}\oplus\eta^{-2}.
\]
Riemann--Hurwitz gives $g(C)=4$, while the cyclic-cover canonical bundle
formula gives
\[
 K_C=\pi^*(K_E\otimes\eta^2)=\pi^*\eta^2=\kappa^2,
 \qquad \kappa=\pi^*\eta.
\]
Projection formula yields
\[
 h^0(C,\kappa)
 =h^0(E,\eta)+h^0(E,\mathcal O_E)+h^0(E,\eta^{-1})=2.
\]

The pullback $\pi^*:J(E)\to J(C)$ is injective.  Indeed, if
$\alpha\in\operatorname{Pic}^0(E)$ and $\pi^*\alpha\simeq\mathcal O_C$, then
\[
 1=h^0(C,\pi^*\alpha)
  =h^0(E,\alpha)+h^0(E,\alpha\eta^{-1})
   +h^0(E,\alpha\eta^{-2}),
\]
and the last two terms vanish; hence $\alpha=\mathcal O_E$.  Duality now
shows that $\ker(\Nm_\pi)$ is connected.  We therefore write
\[
 P=\ker(\Nm_\pi).
\]
It has dimension three.

Use the theta characteristic $\kappa$ to centre the theta divisor on
$J(C)=\operatorname{Pic}^0(C)$:
\[
 \Theta_\kappa
 =\{M\in J(C):h^0(C,\kappa\otimes M)>0\}.
\]
Let $L=\mathcal O_{J(C)}(\Theta_\kappa)|_P$.  The standard polarization
formula for a degree-three cyclic cover ramified at three points gives
\[
 \operatorname{type}(L)=(1,1,3),
 \qquad L^3=3!\,\chi(L)=18;
\]
see \cite[Section~4]{LangeOrtega2011} and
\cite[Corollary~12.1.4 and Lemma~12.3.1]{BirkenhakeLange2004}.
Scheme-theoretically,
\[
 X=\Theta_\kappa\cap P
  =\{M\in P:h^0(C,\kappa\otimes M)>0\}.
\]
Thus $X\in|L|$ and is ample.
Serre duality and $\kappa^2=K_C$ give, for $M\in P$,
\[
 h^0(C,\kappa\otimes M)=h^0(C,\kappa\otimes M^{-1}),
\]
and therefore
\begin{equation}
 X=-X
\label{eq:symmetry}
\end{equation}
with centre the origin.  The uniqueness of this centre will follow from
\eqref{eq:stabilizer}.

\subsection{The norm fibre and the resolution}

Let
\[
 n:C^{(3)}\longrightarrow\operatorname{Pic}^3(E),
 \qquad D\longmapsto\mathcal O_E(\pi_*D),
\]
and put
\[
 N=n^{-1}(\eta^3),
 \qquad
 a:N\longrightarrow X,
 \quad D\longmapsto\mathcal O_C(D)\otimes\kappa^{-1}.
\]
The pencil $\Gamma=|\kappa|\simeq\PP^1$ will be shown below to be the only
positive-dimensional fibre of $a$.  We record the precise resolution because
its topology determines both the Euler characteristic and the Hodge vector.

Let $r_1,r_2,r_3$ be the ramification points and
$D_0=r_1+r_2+r_3\in\Gamma$.  In local coordinates at the three branch
points, the norm map is analytically
\[
 (x_1,x_2,x_3)\longmapsto x_1^3+x_2^3+x_3^3.
\]
Thus $(N,D_0)$ is the affine cone over a smooth plane cubic, a simple
elliptic singularity of type $\wt E_6$ in the standard classification
\cite{Laufer1977,Pinkham1977}.

We next check that there are no other singularities.  Away from the
ramification points the differential of the norm map is nonzero.  At a
ramification point, the local covering is $z\mapsto z^3$.  For a divisor with
two points colliding there, in elementary symmetric coordinates one has
\[
 z_1^3+z_2^3=e_1^3-3e_1e_2,
\]
which has no linear term; with three points colliding, however,
\[
 z_1^3+z_2^3+z_3^3=e_1^3-3e_1e_2+3e_3
\]
has a nonzero linear term.  Hence the only possible critical divisors in
$N$ are $D_0$ and $2r_i+r_j$ with $i\ne j$.  The latter would require
$2b_i+b_j\sim3O$.  Since the $b_i$ are the nonzero two-torsion points, this
would give $b_j=O$, a contradiction.  Thus $D_0$ is the unique singular
point of $N$.

Blowing up its vertex now gives a smooth surface $T$ with exceptional
elliptic curve $E_0$ satisfying $E_0^2=-3$.  The strict transform
$\wt\Gamma$ meets $E_0$ transversely once.

The symmetric-product canonical formula
\cite{Macdonald1962} gives $K_{C^{(3)}}=a_C^*\Theta_C$, and adjunction on the
norm fibre gives $K_N\cdot\Gamma=0$.  Since the singularity has multiplicity
three,
\[
 K_T=\sigma^*K_N-E_0.
\]
Adjunction on $\wt\Gamma$ then yields $\wt\Gamma^2=-1$.  Contracting
$\wt\Gamma$ produces a smooth surface $S$; the image $E_{\exc}$ of $E_0$
satisfies
\[
 E_{\exc}^2=-2.
\]
Finally, $S\to X$ contracts $E_{\exc}$ and is the minimal resolution.  The
standard classification of simple elliptic singularities therefore identifies
the singularity of $X$ at the origin as type $\wt E_7$
\cite{Laufer1977,Pinkham1977}.

For completeness, the Euler characteristic of $N$ can be computed by the
finite map $N\to|\eta^3|\simeq\PP^2$.  Stratifying a line section of the
plane cubic by its multiplicity pattern gives the following table.
\[
\begin{array}{lrrr}
\toprule
\text{base divisor type}&e(\text{stratum})&\#\text{ fibre}&\text{contribution}\\
\midrule
\text{three distinct points, no branch point}&14&27&378\\
\text{three distinct points, one branch point}&-12&9&-108\\
\text{the branch divisor}&1&1&1\\
2q+r,\ q,r\notin B&-24&18&-432\\
2b_i+r&3&3&9\\
2q+b_i&12&6&72\\
3q\ \text{at a flex}&9&10&90\\
\midrule
&&&10
\end{array}
\]
The tangent incidences in this table are distinct: a line through a branch
point $b_i$ and tangent at $q$ satisfies $2q=-b_i$, and multiplication by
two on $E$ is \'etale of degree four.  No nonzero two-torsion point is a flex.
Together with the critical-point calculation above, this proves
\begin{equation}
 e(N)=10,\qquad e(T)=9,\qquad e(S)=8.
\label{eq:euler-chain}
\end{equation}

The canonical map of $C$ is birational: on the affine chart it is
\[
 (x,w)\longmapsto[w^2:w:1:x].
\]
Hence $C$ is not hyperelliptic.  Its canonical quadric is the rank-three
quadric determined by the pencil $|\kappa|$, which has only one ruling.
Consequently $\kappa$ is the unique $g^1_3$ on $C$.  Riemann--Kempf
\cite{Kempf1973} shows that the Abel map is an isomorphism over the $h^0=1$
locus.  The corresponding points of $N$ are smooth by the critical-point
calculation above, so their images in $X$ are smooth.  The pencil
$|\kappa|$ is the only positive-dimensional Abel fibre.  It follows that
\begin{equation}
 \Sing(X)=\{0\}.
\label{eq:singular-locus}
\end{equation}
The smooth dense open $X\setminus\{0\}$ shows that the Cartier divisor $X$
is generically reduced; since a Cartier divisor on the smooth variety $P$ is
Cohen--Macaulay, $X$ is reduced.  An ample divisor on an abelian variety is
connected.  If $X$ were reducible, two components would meet in a curve,
which would lie in $\Sing(X)$.  Thus $X$ is integral.

The map $\rho:S\to X$ is semismall, with the single relevant exceptional
fibre $E_{\exc}$.  Its one-dimensional top intersection form is the
nondegenerate matrix $(-2)$.  The semismall decomposition theorem
\cite{deCataldoMigliorini2002} gives
\begin{equation}
 R\rho_*\QQ_S[2]\simeq\IC_X\oplus\delta_0.
\label{eq:semismall}
\end{equation}
Taking Euler characteristics in \eqref{eq:semismall} and using
\eqref{eq:euler-chain}, we obtain
\begin{equation}
 \chi(\IC_X)=7.
\label{eq:chi-seven}
\end{equation}

The unique singular point also gives
\begin{equation}
 \Stab_P(X)=\{0\},
\label{eq:stabilizer}
\end{equation}
because a translation preserving $X$ must preserve its singular locus.
In particular, the origin is the unique centre of the symmetry
\eqref{eq:symmetry}: two distinct centres would compose to a nontrivial
translation stabilizing $X$.

We will also use dominance of the conormal Gauss map.  Here is a local-index
calculation of its degree.  The normal form in the simple elliptic
classification shows that a general hyperplane section of the $\wt E_7$
germ is of type $A_3$.  Its sectional Milnor number is therefore $3$, and
the local Euler obstruction is $1-3=-2$
\cite{Laufer1977,Pinkham1977,BrasseletLeSeade2000}.  On the other hand,
\eqref{eq:semismall} gives $\chi((\IC_X)_0)=-1$.  The
Dubson--Kashiwara local index formula
\cite{KashiwaraSchapira1990} therefore gives
\[
 CC(\IC_X)=\Lambda_X+\Lambda_0.
\]
The Franecki--Kapranov index theorem \cite{FraneckiKapranov2000}, together
with \eqref{eq:chi-seven}, yields
\begin{equation}
 \deg(\Lambda_X\to\PP(T_0^*P))=6.
\label{eq:gauss-degree}
\end{equation}

\subsection{The generic-twist Hodge vector}

We next determine the Hodge cocharacter needed in the group-theoretic step.
The canonical divisor of $N$ is the pullback of $L$, so $K_N^2=L^3=18$.
The blowup of the $\wt E_6$ vertex gives
\[
 K_T^2=K_N^2+E_0^2=15,
\]
and $T\to S$ is the blowup of a smooth point.  Hence
\begin{equation}
 K_S^2=16,\qquad e(S)=8,\qquad \chi(\mathcal O_S)=2,
\label{eq:surface-invariants}
\end{equation}
where the last equality is Noether's formula.

Intersection-cohomology weak Lefschetz for the ample divisor $X\subset P$
\cite{BBD1982}, together with \eqref{eq:semismall}, identifies
\[
 H^1(P,\QQ)\xrightarrow{\sim}IH^1(X,\QQ)
 \xrightarrow{\sim}H^1(S,\QQ).
\]
Thus $q(S)=3$ and $\operatorname{Pic}^0(P)\to\operatorname{Pic}^0(S)$ is an
isogeny.  The morphism $S\to P$ has a two-dimensional image, so $S$ has
maximal Albanese dimension.

Choose a general unitary rank-one local system $\mathcal L_\alpha$ on $P$,
with associated line bundle $\alpha\in\operatorname{Pic}^0(P)$.  Coherent
generic vanishing \cite{GreenLazarsfeld1987} gives
\[
 H^i(S,K_S\otimes\alpha|_S)=0\qquad(i>0).
\]
Since $\chi(\mathcal O_S)=2$, Riemann--Roch and Serre duality, applied also
to $\alpha^{-1}$, give
\[
 h^0(S,K_S\otimes\alpha|_S)=2,
 \qquad h^2(S,\alpha|_S)=2.
\]
Let $j:S\to P$ be the proper morphism used below.  The complex
$Rj_*\QQ_S[2]$ is perverse: first $R\rho_*\QQ_S[2]$ is perverse because
$\rho$ is semismall, and then extension along $X\hookrightarrow P$ is
perverse $t$-exact.  Generic vanishing for perverse sheaves on abelian
varieties \cite{KraemerWeissauer2015,Schnell2015} therefore gives
\[
 H^k(S,j^*\mathcal L_\alpha)=0\qquad(k\ne2).
\]
For a unitary local system, the Hodge decomposition is the twisted Dolbeault
decomposition
\[
 H^{p,q}(S,j^*\mathcal L_\alpha)
 =H^q(S,\Omega_S^p\otimes\alpha|_S),
 \qquad p+q=2.
\]
The total dimension is $e(S)=8$; the two extreme dimensions are both $2$,
so the Hodge vector of $Rj_*\QQ_S^H[2]$ is $(2,4,2)$.  The Hodge-module
refinement of \eqref{eq:semismall}, extended to $P$, is
\[
 Rj_*\QQ_S^H[2]\simeq \IC_X^H\oplus\QQ_0^H(-1).
\]
The point summand has type $(1,1)$ and removes one middle Tate class.  We
have proved:

\begin{proposition}\label[proposition]{prop:hodge-vector}
The generic-twist Hodge vector of $\IC_X^H$ is
\[
 h(\IC_X^H)=(2,3,2).
\]
The associated antidiagonal Hodge cocharacter has weights
$-2,0,2$ on $\omega(\IC_X)$ with multiplicities $2,3,2$.
\end{proposition}

\section{The universal pair-incidence splitting}\label{sec:pair-splitting}

We now prove the geometric input for pair convolution.  All statements in
this section are made after restricting to the open loci on which the
displayed divisors are reduced and all Abel points have $h^0=1$.

\subsection{A dominant normalized chart}

Let $D$ be the unique reduced degree-three divisor representing a general
point $x\in X$.  Its norm is a general line section of the plane cubic $E$.
After an affine coordinate change, the line can be written $y=x$.  If the
three selected lifts have $w$-coordinates $a,b,c$, then
\[
 f(x)-x^2=\lambda(x-a^3)(x-b^3)(x-c^3),\qquad \lambda\ne0.
\]
Put $e_i=e_i(a,b,c)$ and
\[
 p(W)=W^3-e_1W^2+e_2W-e_3.
\]
Conversely, these formulas reconstruct the curve and the divisor
\[
 D:\quad x=y=w^3,\qquad p(w)=0.
\]
Thus $(\lambda,e_1,e_2,e_3)$ is a dominant rational chart of the actual
space of pairs $(C,D)$.

There is a scaling redundancy
\[
 (x,y,w)\longmapsto(r^3x,r^3y,rw),
\]
under which
\[
 (\lambda,e_1,e_2,e_3)
 \longmapsto(r^{-3}\lambda,re_1,r^2e_2,r^3e_3).
\]
On $e_1\ne0$ the rational choice $r=e_1^{-1}$ gives the gauge $e_1=1$.
No root is adjoined, so this normalization does not split a quadratic cover
artificially.

\subsection{The norm conic}

The evaluation map on $D$ has rank three.  Hence
$H^0(C,3\kappa-D)$ has dimension three, with basis
\[
 s_0=x-y,
 \qquad
 s_1=p(w),
 \qquad
 s_2=(x-e_3)w-e_1y+e_2w^2.
\]
For $s=us_0+vs_1+zs_2$, write
\[
 s=\alpha+\beta w+\gamma w^2,
\]
where
\[
\begin{aligned}
 \alpha&=u(x-y)+v(y-e_3)-e_1zy,\\
 \beta&=e_2v+z(x-e_3),\\
 \gamma&=-e_1v+e_2z.
\end{aligned}
\]
Since $w^3=y$, its cyclic norm is
\begin{equation}
 \Nm(s)=\alpha^3+\beta^3y+\gamma^3y^2-3\alpha\beta\gamma y.
\label{eq:cyclic-norm}
\end{equation}
The section \eqref{eq:cyclic-norm} is divisible on $E$ by the fixed line
$y-x$.  The quotient belongs to $H^0(E,6O)$ and is represented by a conic
\[
 Q_s=q_{00}+q_{01}x+q_{02}y+q_{11}x^2+q_{12}xy+q_{22}y^2.
\]
Let $M_s$ be its symmetric $3\times3$ matrix.

\begin{proposition}[Universal norm identity]
\label[proposition]{prop:universal-identity}
Over $\QQ(\lambda,e_1,e_2,e_3)$ one has
\[
 \det(M_s)=L_{\mathrm{str}}F_8,
\]
where $L_{\mathrm{str}}=u-v$ has multiplicity one and $F_8$ is a homogeneous
factor of degree eight.  The geometric generic curve $F_8=0$ is integral.

On the chart $e_1=v=1$, let $F(z)$ be the monic sextic obtained from $F_8$
and put
\[
 \Delta=-4\left(q_{00}q_{11}-(q_{01}/2)^2\right).
\]
There exist polynomials $A,B\in\QQ(\lambda,e_2,e_3,u)[z]$ of degrees three
and at most two such that
\begin{equation}
 F=A^2+3\Delta B^2.
\label{eq:norm-square-identity}
\end{equation}
In particular, the intrinsic double cover ordering the two factors of a
rank-two conic splits over the function field of $F_8$ after base change to
$\CC$.
\end{proposition}

\begin{proof}
The determinant factorization follows by substituting
\eqref{eq:cyclic-norm}, reducing modulo $y^2=f(x)$, and solving in the basis
$1,x,y,x^2,xy,y^2$ of $H^0(E,6O)$.  All operations are over the indicated
rational function field.

The residual factor is geometrically integral on a nonempty open.  At the
specialization
\[
 (\lambda,e_1,e_2,e_3)=(1,6,11,6),
 \qquad (a,b,c)=(1,2,3),
\]
its reduction modulo $13$ is irreducible of degree eight
and has the smooth rational point $(u,v)=(0,2)$ on the chart used in the
certificate.  If an irreducible curve over a finite field were geometrically
reducible, Frobenius would permute its geometric components transitively;
every rational point would lie on every component and would be singular.
Thus this fibre is geometrically integral.  Openness of geometric
integrality in a flat projective family proves the generic assertion.

The exact coefficient formulas and two independent checks of
\eqref{eq:norm-square-identity} are given in \cref{app:certificate}.  On the
cover $t^2+3\Delta=0$, the identity becomes
\[
 F=(A+tB)(A-tB).
\]
The polynomial $B$ is nonzero and has degree at most two, whereas $F$ is
irreducible of degree six on this chart; hence $B$ is nonzero in the
function field.  The two rational roots $t=\pm A/B$ are distinct.  Finally,
$\Delta$ is the discriminant of the restriction of $Q_s$ to a coordinate
line.  On the rank-two locus, ratios obtained from different nonzero
principal minors are squares because $\operatorname{adj}(M_s)$ has rank
one.  Thus $\Delta$ defines the intrinsic ordering character.  The constant
$-3$ is a square over $\CC$, completing the proof.
\end{proof}

\subsection{The scheme-theoretic Abel inverse}

We isolate the relative argument that turns ordered conic factors into
ordered components of the Abel incidence.

\begin{lemma}[Relative ordered-zero construction]
\label[lemma]{lem:scheme-abel-inverse}
Let $S$ be a nonempty open of the unordered distinct rank-two incidence.
Let $\mathcal I_{\mathrm{fac}}\to S$ be the finite \'etale double cover that
orders the two conic factors.  Let $\mathcal I_{\mathrm{div}}\to S$ be the
incidence whose $T$-points are ordered relative degree-three Cartier
divisors $(D_1,D_2)$ with
\[
 \mathcal O_{C_T}(D+D_1+D_2)\simeq3\kappa_T,
 \qquad \pi_*D_i\in|\eta_T^3|,
\]
such that the induced projective section of $3\kappa_T$ is the pullback of
$[s]$ and the two norm divisors are the residual line sections of $Q_s$.
Impose also the transverse, disjointness, branch-avoidance, and Abel-open
conditions used below.  After further shrinking $S$, the following
assignments define inverse $S$-isomorphisms
$\mathcal I_{\mathrm{fac}}\simeq\mathcal I_{\mathrm{div}}$ and commute with
every base change $S'\to S$, including nonreduced base change:
\begin{enumerate}[label=\textup{(\roman*)}]
\item factor the residual norm conic as an ordered product
      $Q_s=\ell_1\ell_2$;
\item write
      $\operatorname{div}(s)=D+D_1+D_2$, where
      $\Nm(D_i)$ is the line section cut out by $\ell_i$.
\end{enumerate}
The divisors $D_i$ are relative effective Cartier divisors of degree three,
and, after restricting to the relative smooth Abel open, their Abel classes
are smooth points of $X$.
\end{lemma}

\begin{proof}
For the forward construction put $T=\mathcal I_{\mathrm{fac}}$ and write
$Q_s=\ell_1\ell_2$ in its universal ordering.  Remove the loci where a line
is tangent to $E$, where either line meets the
branch divisor, where any two of the fixed and residual line sections meet,
where $s$ vanishes on two sheets over one point of a residual line section,
where a norm factor has multiplicity greater than one, or where a relevant
Abel class has $h^0>1$ or lies outside the relative smooth locus of $X/T$.
The line $\ell_i$ then
cuts out a relative finite \'etale divisor $B_i\subset E_T$ of degree three.
Moreover
\[
 U_i=C_T\times_{E_T}B_i\longrightarrow B_i
\]
is finite \'etale of degree three, because $B_i$ avoids the branch divisor.

We first select the zero over $B_i$ scheme-theoretically.  \'Etale-locally on
$B_i$, split $U_i=B_i\sqcup B_i\sqcup B_i$ and trivialize the restriction
of $3\kappa$, using the induced trivialization of its norm line.  If
$a_1,a_2,a_3\in\mathcal O_{B_i}$ are the three values of $s$, compatibility
of the finite locally free norm with base change gives
\begin{equation}
 a_1a_2a_3=\Nm_\pi(s)|_{B_i}=0
\label{eq:split-norm-zero}
\end{equation}
as an equality in the structure sheaf, not merely at geometric points.
The collision exclusions imply that locally two values, say $a_2,a_3$, are
units.  Equation~\eqref{eq:split-norm-zero} then forces $a_1=0$, also in
nilpotent directions.
Consequently
\[
 Z_i=Z(s)\times_{E_T}B_i\subset U_i
\]
is locally exactly one open-and-closed sheet of $U_i/B_i$.  Thus
$Z_i\to B_i$ is finite \'etale of degree one and hence is an isomorphism.

Let $t$ be a local equation for $B_i$ in $E_T$.  Since $B_i$ occurs with
multiplicity one in the norm divisor and the other norm factors are
disjoint from it, $\Nm_\pi(s)=t u$ with $u$ a unit along $B_i$.  \'Etale-locally
on a neighbourhood of $B_i$ in the complement of the branch divisor, split
the cover $C_T\to E_T$.  On the selected sheet the other two values of $s$
remain units, so the same local product formula shows that the selected
value is $t$ times a unit.
The zero is therefore simple.  The image $D_i\subset C_T$ of $Z_i$ is
finite \'etale of degree three over $T$; as a finite \'etale multisection of the
smooth relative curve $C_T/T$, it is a relative effective Cartier divisor.

The pairwise disjoint divisors $D,D_1,D_2$ lie in $Z(s)$ and have total
relative degree nine, equal to $\deg(3\kappa)$.  After dividing by their
Cartier equations, the residual section has degree zero and is nonvanishing
on every geometric fibre.  Nakayama's lemma makes it a trivialization, so
there is no vertical or nilpotent residual contribution and
\begin{equation}
 Z(s)=D+D_1+D_2
\label{eq:relative-zero-sum}
\end{equation}
as relative Cartier divisors.  Conversely, an ordered divisor datum
determines the unique projective section $[s]$ with zero divisor
$D+D_1+D_2$, and taking its norm recovers the ordered line sections $B_i$
and hence the ordered factors $\ell_i$.  Thus these assignments are inverse
already at the divisor level.

For completeness, let
\[
 \alpha:C_T^{(3)}\longrightarrow\operatorname{Pic}^3(C_T/T)
\]
be the relative Abel map and let
$F=\Nm^{-1}(\eta^3)\subset\operatorname{Pic}^3(C_T/T)$.  The norm fibre is
the cartesian base change
\[
 N_T=C_T^{(3)}\times_{\operatorname{Pic}^3(C_T/T)}F.
\]
On the relative Picard scheme, let $W_3^{(1)}$ be the open on which $h^0=1$
and the pushforward of a Poincar\'e bundle is locally free of rank one and
commutes with base change; pull this open back to $T$.  The standard
projective-bundle description of the Abel map and cohomology and base change
then identify its source with the projectivization of that rank-one bundle.
Thus the map is an isomorphism and remains so after arbitrary base change.
Fibrewise, this is the $h^0=1$ case of Riemann--Kempf
\cite{Kempf1973}:
\[
 \alpha^{-1}(W_3^{(1)})\xrightarrow{\sim}W_3^{(1)}.
\]
Base changing this isomorphism by $F\to\operatorname{Pic}^3(C_T/T)$ and
translating by $\kappa^{-1}$ gives
\[
 N_T^{(1)}\xrightarrow{\sim}X_T^{(1)}.
\]
We replace both sides by the inverse images of
$X_T^{(1)}\cap X_{T,\mathrm{sm}/T}$, as imposed above.  Thus the divisor
construction descends to the relative smooth Abel incidence.

Zero schemes, fibre products, finite locally free norms, and the cartesian
Abel construction all commute with arbitrary base change.  Hence both
composites preserve the universal zero schemes and are the identity as
morphisms.  The use of $C_T^{(3)}\simeq\operatorname{Hilb}^3(C_T/T)$ is only
representability; no normalization of the ordering cover or of pair fibres
is involved.
\end{proof}

The open in \cref{lem:scheme-abel-inverse} is nonempty.  After reduction
modulo $13$, an exact witness is provided by
\[
 E:y^2=x^3-64x+64,
 \quad (a,b,c)=(1,2,-2),
 \quad (u:v:z)=(1:0:9)\pmod {13}.
\]
At this point the conic factors over $\mathbf F_{13}$ as
\[
 -(4x+y-6)(6x+y+5).
\]
The two line-section discriminants, the branch resultants, the pair
resultant, and the six evaluations on the fixed divisor are all nonzero;
see \cref{app:certificate}.  The unique $g^1_3$ is $|\kappa|$.  In the
coordinates above, the norm of a section of
$H^0(C,\kappa)=H^0(E,\eta)\oplus H^0(E,\mathcal O_E)$ cuts out a line with
no $x$-term, whereas the three norm lines
\[
 y-x,\qquad 4x+y-6,\qquad 6x+y+5
\]
all have a nonzero $x$-term.  Thus their divisors have $h^0=1$ and avoid the
contracted pencil.  Hence every excluded condition above is a proper closed
condition.

\subsection{Boundary exhaustion and global labels}

Fix a geometric generic pair $(C,D_x)$.  For
\[
 C_x=\{y\in X:-x-y\in X\},
\]
the unique effective divisors representing $x,y,-x-y$ define a rational map
\[
 \Psi:C_x\dashrightarrow\PP H^0(C,3\kappa-D_x).
\]
Its fibres are finite: for a fixed section, the residual zero-cycle has only
finitely many degree-three subcycles.  Thus every curve component of $C_x$
maps to a component of $\det(M_s)=0$.

The structural line has the intrinsic description
\[
 \PP H^0(3\kappa-D_x-D_{-x})=\PP H^0(\kappa).
\]
Its members have divisor $D_x+D_{-x}+D_0$ with
$D_0\in|\kappa|$.  For two reduced disjoint line sections of a plane cubic,
a mixed degree-three subcycle is not a line section: the line through two
points of either triple contains the third point of that same triple.  Hence
the structural line represents only the fixed pair $\{-x,0\}$ and does not
give a curve component.  The indeterminacy locus $y=0$ or $-x-y=0$ is also
zero-dimensional.  Since $F_8$ is the only remaining determinant component,
\cref{lem:scheme-abel-inverse} exhausts all nonexceptional curve components.
We conclude that the geometric generic $C_x$ is reduced and has exactly two
integral components $A_x,B_x$, exchanged by
\[
 y\longmapsto-x-y.
\]
Reducedness follows because $C_x$ is a local complete intersection curve in
the smooth threefold $P$, hence Cohen--Macaulay, and is reduced at both
generic points.

It remains to descend the two labels from the normalized chart.  Let $\tau$
be the order-three deck transformation.  On $P$ one has
\[
 1+\tau+\tau^2=0,
\]
and $\tau(X)=X$.  Therefore
\[
 s_+(x)=(\tau x,\tau^2x)\in C_x
\]
is a rational section of the pair-incidence family.  It is smooth for
general $x$.  Otherwise the Gauss covectors at $\tau x$ and $\tau^2x$ would
agree; equivariance would force the general Gauss image into the projective
fixed locus of $\tau$, contrary to \eqref{eq:gauss-degree}.  Explicitly,
\[
 T_0^*P=
 \left\langle\frac{dx}{w^4}\right\rangle
 \oplus
 \left\langle\frac{dx}{w^5},\frac{x\,dx}{w^5}\right\rangle,
\]
and the two summands have distinct $\tau$-eigenvalues.  The projective fixed
locus is a point union a line, a proper subset of $\PP^2$.

Over the function field of the actual pair space, a smooth rational point
lies on exactly one geometric component.  Galois therefore fixes the
component containing $s_+$ and, since there are exactly two components,
also fixes the other.  Taking closures and using generic flatness and
openness of geometric integrality and reducedness gives the following.

\begin{theorem}[Pair splitting]\label{thm:pair-splitting}
After shrinking to a nonempty open subfamily of curves, there is a dense
connected open $U\subset X_{\reg}\setminus\{0\}$ such that, for every
$x\in U$, the nonexceptional part of $C_x$ has exactly two reduced integral
curve components $A_x,B_x$.  They are fibres of two globally labelled flat
relative components over $U$, and factor transposition exchanges them.
\end{theorem}

\section{Pair convolution detects \texorpdfstring{$\IC_X$}{IC(X)}}

Let $\rho:S\to X$ be the minimal resolution and let $j:S\to P$ be its
composition with the inclusion.  Put
\[
 K=\IC_X,
 \qquad
 L=Rj_*\QQ_S[2]=K\oplus\delta_0.
\]
Consider
\[
 \wt m:S\times S\longrightarrow P,
 \qquad (p,q)\longmapsto j(p)+j(q),
\]
and let $\sigma(p,q)=(q,p)$.  We use the minus sign to denote the
anti-invariant direct summand under $\sigma$.

\begin{proposition}[Alternating support detector]
\label[proposition]{prop:alt-support}
Under the conclusion of \cref{thm:pair-splitting}, one has a direct summand
\[
 K\subset{}^pH^0\bigl(Rm_*(K\bo K)\bigr)_-.
\]
Consequently $K$ is a direct summand of $\Alt^2(K)$ in the convolution
Tannakian category.
\end{proposition}

\begin{proof}
Put $U^+=-U$.  Fix $t\in U^+$, write $x=-t\in U$, and let
$E_{\exc}=\rho^{-1}(0)$.  By definition, $C_x$ parametrizes ordered pairs
whose sum is $-x=t$.  Hence the fibre
$F_t=\wt m^{-1}(t)$ has precisely four irreducible curve components: the
strict transforms of $A_x,B_x$ and the two structural components
\[
 C_{1,t}=E_{\exc}\times\{s_t\},
 \qquad
 C_{2,t}=\{s_t\}\times E_{\exc},
\]
where $s_t$ is the unique point above $t$.  Indeed, away from
$E_{\exc}$ the resolution is an isomorphism, while a pair with one entry on
$E_{\exc}$ must have the other entry over $t$.  Thus there are no further
top-dimensional components.  The involution $\sigma$ exchanges $A_x$ with
$B_x$ and $C_{1,t}$ with $C_{2,t}$.

Proper base change identifies
\[
 \cH^{-2}(L*L)_t=H^2(F_t,\QQ).
\]
The top cohomology of a proper complex curve is freely generated by the
fundamental classes of the irreducible components of its reduction.
Holomorphic transposition preserves complex orientation.  The two exchanged
pairs therefore give
\begin{equation}
 \dim\cH^{-2}(L*L)_{-,t}=2.
\label{eq:total-alt-rank}
\end{equation}

Expand the decomposition $L=K\oplus\delta_0$:
\[
 L*L=(K*K)\oplus(K*\delta_0)\oplus(\delta_0*K)
      \oplus(\delta_0*\delta_0).
\]
At $t\ne0$, the last term has zero stalk.  The two middle terms are both
canonically $K$, and transposition exchanges them.  Since $t$ is a smooth
point of the surface $X$, their anti-invariant stalk in degree $-2$ has
dimension one.  Subtracting this from \eqref{eq:total-alt-rank} gives
\begin{equation}
 \dim\cH^{-2}(K*K)_{-,t}=1.
\label{eq:principal-alt-rank}
\end{equation}
There is no Koszul sign: $K$ has geometric shift $[2]$, and
$(-1)^{2\cdot2}=1$.

After a topological stratification and a further shrinking of $U^+$, the four
globally labelled relative components trivialize the full alternating top
cohomology local system.  The anti-diagonal of the two cross terms is also a
constant local subsystem.  Hence the rank-one local system in
\eqref{eq:principal-alt-rank} is constant.

For general $t\in P\setminus X$, the fibre of pair addition is the smooth
connected curve $X\cap(t-X)$.  It is connected because it is the
intersection of two ample divisors in the smooth projective threefold $P$.
The involution $p\mapsto t-p$ preserves its complex orientation and hence
acts by $+1$ on its one-dimensional top cohomology.  Therefore
\begin{equation}
 \cH^{-2}(K*K)_-=0
 \quad\text{on a dense open subset of }P.
\label{eq:generic-alt-vanishing}
\end{equation}

The complex $(K*K)_-$ is a direct summand of the projective direct image of
a pure polarizable Hodge module.  By the decomposition theorem and
semisimplicity \cite{BBD1982,Saito1988}, it is a direct sum of shifted simple
intersection complexes.  A summand with proper support containing the
generic point of $U^+$ must have support $X$.  Since
$\IC_X(\mathcal V)|_{U^+}=\mathcal V[2]$, contribution to ordinary degree
$-2$ forces perverse shift zero.

We spell out why a full-support summand cannot account for the line in
\eqref{eq:principal-alt-rank}.  At a general point of $U^+$ choose a normal
analytic disk $\Delta$ to $X$, and write $j_\Delta:\Delta^*\hookrightarrow
\Delta$.  Locally, a simple perverse intersection complex with support $P$
has the form
\[
\mathcal L[2]\boxtimes (j_{\Delta,!*}\mathcal W[1])
\]
for local systems $\mathcal L$ and $\mathcal W$.  Writing
$i_0:\{0\}\hookrightarrow\Delta$, the defining no-quotient property of
intermediate extension gives
\[
 \cH^0\bigl(i_0^*j_{\Delta,!*}\mathcal W[1]\bigr)=0.
\]
Since a perverse sheaf on a disk has ordinary stalk amplitude $[-1,0]$,
$i_0^*j_{\Delta,!*}\mathcal W[1]$ is concentrated in degree $-1$ when
nonzero.  Thus an unshifted
full-support summand can contribute there only in ordinary degree $-3$.
Contribution in degree $-2$ requires the perverse-degree-one shift $[-1]$.
But on the dense smooth open of $P$ that same summand is a nonzero local
system shifted by $[3][-1]$ and therefore contributes in degree $-2$,
contrary to \eqref{eq:generic-alt-vanishing}.  Supports of dimension at most
one do not meet a generic point of $U^+$, and a different irreducible divisor
cannot contain $U^+$.

Thus the constant rank-one local system in
\eqref{eq:principal-alt-rank} is the coefficient system of an unshifted
$\IC_X$ summand.  This proves the first assertion.  The even support
dimension identifies the anti-invariant convolution summand with the
ordinary exterior square under the fibre functor.  Finally, generic
vanishing makes the nonzero perverse-cohomological degrees of convolution
negligible; the semisimple quotient is exact and preserves direct summands.
Since $\chi(K)=7$, the displayed $K$ remains nonzero in the quotient.  This
proves the second assertion.
\end{proof}

\section{Identification of the full Tannaka group}

Let
\[
 G=G(K),\qquad V=\omega(K).
\]
By \eqref{eq:chi-seven}, $\dim V=7$.  The symmetry \eqref{eq:symmetry},
Verdier self-duality, and the even support dimension give a perfect symmetric
evaluation pairing
\[
 K*K\longrightarrow\delta_0.
\]
It takes values in the actual tensor unit because the centre of the symmetry
is zero.  Therefore the full group embeds faithfully as
\begin{equation}
 G\hookrightarrow O(V).
\label{eq:orthogonal}
\end{equation}

The simple perverse sheaf $K$ is nondivisible: a translation fixing it must
fix its support, and \eqref{eq:stabilizer} then makes the translation zero.
By \cite[Proposition~3.3]{JKLM2025},
\begin{equation}
 V|_{G^0}\quad\text{is irreducible}.
\label{eq:connected-irreducible}
\end{equation}
This result does not require the ambient abelian variety to be simple.

Let $G_P=G$ and let $G_H$ be the Tannaka group of the pure Hodge module
$\IC_X^H$, acting faithfully on the same complex vector space $V$.  By
\cite[Corollary~3.2]{KLM2026}, the Hodge decomposition is defined by a
cocharacter $\lambda:\GGm^2\to G_H$.  The comparison theorem
\cite[Corollary~2.5]{KLM2026} gives
\[
 G_P\mathrel{\triangleleft}G_H,
 \qquad
 [G_P^0,G_P^0]=[G_H^0,G_H^0]=:G^*.
\]
Since $G_P^0\subset G_H^0$, equation
\eqref{eq:connected-irreducible} implies that $V|_{G_H^0}$ is irreducible.
Thus the connected centre $Z_H^0$ acts by scalar matrices.  As $G_H^0$ is
reductive, multiplication $Z_H^0\times G^*\to G_H^0$ is surjective with
finite kernel.

Put $\rho(z)=\lambda(z,z^{-1})$.  The source of $\lambda$ is connected, so
$\rho(\GGm)\subset G_H^0$.  By \cref{prop:hodge-vector}, $\det\rho=1$.
The determinant is trivial on the semisimple group $G^*$ and therefore
descends to a character
\[
 \overline{\det}:G_H^0/G^*\longrightarrow\GGm.
\]
The quotient $G_H^0/G^*$ is the image of $Z_H^0$.  On $Z_H^0$ the
determinant is $c\mapsto c^7$; since the faithful scalar action embeds
$Z_H^0$ into $\GGm$, the induced character $\overline{\det}$ has finite
kernel.  Hence the image of the connected group $\GGm$ under $\rho$ in
$G_H^0/G^*$ is both connected and finite, and is therefore trivial.
Consequently
\[
 \rho(\GGm)\subset G^*,
\]
the common Hodge and perverse derived group.  Its weights on $V$ are
\begin{equation}
 -2,0,2\quad\text{with multiplicities }2,3,2.
\label{eq:hodge-weights}
\end{equation}

Applying the fibre functor to \cref{prop:alt-support} gives
\begin{equation}
 \operatorname{Hom}_G(V,\Lambda^2V)\ne0.
\label{eq:hom-exterior}
\end{equation}

\begin{proposition}\label[proposition]{prop:connected-g2}
The identity component of $G$ is $G_2$ in its standard representation.
\end{proposition}

\begin{proof}
By \eqref{eq:orthogonal} and \eqref{eq:connected-irreducible}, the connected
centre of $G^0$ acts by orthogonal scalars and is therefore trivial.  Thus
$G^0$ is semisimple.  An irreducible representation of a product of
semisimple groups is an external tensor product.  Since $7$ is prime and the
action is faithful, $G^0$ has one simple factor.

The classification of irreducible self-dual seven-dimensional
representations leaves exactly
\[
\begin{gathered}
 SO_7\ \text{on its standard module},\qquad
 G_2\ \text{on its standard module},\\
 PGL_2\ \text{on }\Sym^6(\CC^2).
\end{gathered}
\]
Passing to the simply connected cover and consulting the low-dimensional
highest weights gives this list directly
\cite{Dynkin1957,FultonHarris1991}: type $A_6$ also has a seven-dimensional
standard module, but it is not self-dual, and the remaining simple types
have no irreducible module of dimension seven.  For $SO_7$, the exterior
square of the
standard module is the irreducible adjoint module, so
$\operatorname{Hom}_{SO_7}(V,\Lambda^2V)=0$, contrary to
\eqref{eq:hom-exterior}.  For a nontrivial cocharacter of $PGL_2$, the
weights on $\Sym^6(\CC^2)$ are seven distinct numbers.  For a primitive
cocharacter they are
\[
 -3,-2,-1,0,1,2,3,
\]
and every other nonzero cocharacter multiplies them by a nonzero integer.
The trivial cocharacter has one weight of multiplicity seven.  Neither
possibility agrees with \eqref{eq:hodge-weights}.  Hence $G^0=G_2$.
\end{proof}

It remains to identify the component group.  In the standard orthogonal
representation,
\[
 N_{O_7}(G_2)=G_2\times\{\pm I\}.
\]
Indeed, $G_2$ has no outer automorphisms, and its centralizer on the
irreducible standard module consists of scalars.  If $G$ met the second
coset, it would contain $-I$.  But $-I$ acts by $-1$ on $V$ and by $+1$ on
$\Lambda^2V$, forcing every morphism in \eqref{eq:hom-exterior} to vanish.
This contradiction proves $G=G_2$.

\section{The nonempty open family}\label{sec:open-family}

We now make the quantifier in \cref{thm:main} explicit.  Over the regular
integral base $\cB$, let $\cC\to\cB$ be the relative family of cyclic covers,
let $\cP\to\cB$ be the relative Prym abelian scheme, and let
$\cX\subset\cP$ be the relative centred theta section.  The divisor
$\cX$ is flat over $\cB$: it is a relative effective Cartier divisor in the
smooth scheme $\cP$, and its defining section remains a non-zero-divisor on
every fibre.  The preceding geometry applies after arbitrary geometric base
change: every fibre is geometrically integral and its only singular point is
the zero section.  Put
\[
 \cM=(\cX/\cB)_{\mathrm{sm}}=\cX\setminus 0.
\]
Then $\cM\to\cB$ is smooth and surjective, and every fibre is a dense
geometrically integral open of the corresponding theta surface.  Moreover,
$\cM$ is integral.  Indeed, since $\cB$ is regular and $\cM\to\cB$ is
smooth, $\cM$ is regular, so its irreducible components are open.  The image
of each component is a nonempty open of the integral base $\cB$ and hence
contains its generic point.  Each component would therefore yield an
irreducible component of the geometrically integral generic fibre.  Thus
there is only one.

Let $\cN$ be the relative norm fibre in $\cC^{(3)}$.  On $\cM$ every Abel
class has $h^0=1$, so the relative Abel isomorphism used in
\cref{lem:scheme-abel-inverse} identifies $\cN\times_{\cX}\cM$ with $\cM$.
The universal divisor on the relative symmetric cube therefore supplies the
universal effective divisor over $\cM$.  The unnormalized four-parameter
chart $(\lambda,e_1,e_2,e_3)$ of \cref{sec:pair-splitting}, before imposing
the gauge $e_1=1$, maps dominantly to $\cM$.  After restricting its domain,
let $T\to\cM$ be an integral dominant chart.  Equivalently, one may
reintroduce the scaling parameter $r\in\GGm$ into the slice $e_1=1$; that
slice is used only for the function-field calculation and does not lower the
dimension of the dominant chart.  After the standard local-freeness
shrinking, write $p:\cC_T\to T$ for the relative curve and $D$ for its
universal degree-three divisor.  Then
$p_*\mathcal O_{\cC_T}(3\kappa-D)$ is a rank-three vector bundle.

All charts, the residual equation, and the excluded-locus functions are
defined over an integral localization
\[
 R_0=\mathbf Z[1/N],
\]
where $13\nmid N$ and every displayed denominator is invertible at the two
modulo-$13$ witnesses.  Denote the resulting integral model of $T$ and its
relative curve by $T_{R_0}$ and
$p_{R_0}:\cC_{T_{R_0}}\to T_{R_0}$.  After shrinking around the relevant
charts, the structure map $T_{R_0}\to\operatorname{Spec}R_0$ is dominant.
Since $R_0$ is a principal ideal domain and $T_{R_0}$ is integral, its affine
coordinate rings are $R_0$-torsion-free; hence $T_{R_0}$ is $R_0$-flat.  The
residual degree-eight Cartier divisor
\[
 \cR_{R_0}\subset
 \PP_{T_{R_0}}\bigl(
 p_{R_0*}\mathcal O_{\cC_{T_{R_0}}}(3\kappa-D)\bigr)
\]
is projective and flat over $T_{R_0}$.  The first modulo-$13$ certificate in
\cref{prop:universal-identity} is a geometrically integral fibre of this
family.  Geometric integrality is open for a flat projective family; the
corresponding nonempty open of the integral scheme $T_{R_0}$ contains its
generic point.  After base change to $\CC$ and a further shrinking of $T$,
it follows that $\cR=\cR_{R_0}\times_{R_0}\CC$ is integral with
geometrically integral fibres.

On this single integral space $\cR$, reducedness of the fixed and residual
degree-three divisors, rank exactly two of the norm conic, distinctness of
its factors, finite \'etaleness and disjointness of the line sections, and
avoidance of the branch divisor and contracted pencil are open conditions.
Let $g$ be the product of the corresponding discriminants, resultants, and
evaluations on the affine witness chart.  The model $\cR_{R_0}$ is flat over
$R_0$ by composition, so its affine coordinate rings are $R_0$-torsion-free.
The second
certificate gives $g\ne0$ at an $\mathbf F_{13}$-point of $\cR_{R_0}$;
hence $g$ is nonzero on the characteristic-zero generic fibre.  Consequently
the good locus $\cR_g\subset\cR$ is a nonempty open.
The universal identity \eqref{eq:norm-square-identity} splits the ordering
cover at the generic point of $\cR$, and the relative ordered-zero
construction identifies its ordered sheets with the two Abel components.
Let $\cI\to\cM$ denote the relative pair-incidence family, whose fibre over
$x$ is $C_x$.  Since $T\to\cM$ is dominant, the preceding calculation after
base change to $T$ shows that the geometric generic fibre of
$\cI\to\cM$ has exactly two geometrically integral reduced components.  The
$\CC(\cM)$-rational section $x\mapsto(\tau x,\tau^2x)$ is a smooth point of
that generic fibre and hence lies on exactly one geometric component.  Galois
fixes this component and, since there are exactly two, also fixes the other.
Thus both components descend to the function field of $\cM$.  Take their
closures $\cA_1,\cA_2\subset\cI$.  By generic flatness and openness of geometric
integrality and reducedness, after replacing $\cM$ by a nonempty open
$\cM^\circ$, the two closures are flat with geometrically integral reduced
curve fibres, transposition exchanges them, and their union is the full
nonexceptional pair fibre.  Shrinking also removes every vertical component.

The morphism $\cM^\circ\to\cB$ is the restriction of a smooth morphism and is
therefore open.  Its image
\[
 \cB^\circ=\operatorname{im}(\cM^\circ\to\cB)
\]
is consequently a nonempty Zariski-open subset.  For every
$b\in\cB^\circ(\CC)$, the fibre
$U_b=\cM^\circ\cap X_b$ is a nonempty open of the integral surface $X_b$;
hence it is dense and connected.  The pair-splitting theorem holds over
$U_b$, and all arguments of the preceding sections apply.  This proves
\cref{thm:main} with the stated ``for every $b$'' quantifier.

\appendix

\section{Exact norm certificate}\label{app:certificate}

We give enough information to reproduce the only large symbolic identity in
the proof.  On $e_1=v=1$, write
\[
 F=z^6+c_5z^5+c_4z^4+c_3z^3+c_2z^2+c_1z+c_0
\]
and put $k=3\Delta$.  Let $L$ be the leading numerator of the unnormalized
sextic obtained from the determinant calculation.  Set
\[
 A=z^3+a_2z^2+a_1z+a_0,
 \qquad B=b_2z^2+b_1z+b_0,
\]
where
\[
\begin{aligned}
b_2=\frac{e_2\lambda}{2L}(&2e_2^3\lambda u+e_2^3\lambda
-6e_2e_3\lambda u+6e_2e_3\lambda-3e_3\lambda+3e_3),\\
b_1=-\frac{\lambda}{2L}(&e_2^3\lambda u+2e_2^3\lambda
+3e_2^2e_3\lambda u^2-6e_2^2e_3\lambda u+3e_2^2e_3\lambda\\
&-3e_2e_3\lambda u+3e_2e_3\lambda-3e_3\lambda+3e_3),\\
b_0=\frac{\lambda(u-1)}{2L}(&-e_2^3\lambda+3e_2e_3\lambda u
-3e_2e_3\lambda+2e_3\lambda-2e_3),
\end{aligned}
\]
and
\[
 a_2=\frac{c_5}{2},\qquad
 a_1=\frac{c_4-a_2^2-kb_2^2}{2},\qquad
 a_0=\frac{c_3}{2}-a_2a_1-kb_2b_1.
\]
Substitution proves $F=A^2+kB^2$.  The verifier performs this substitution
in the exact fraction field $\QQ(\lambda,e_2,e_3,u)$ and obtains seven zero
coefficient remainders.  An independent verifier clears the common
denominator $4\lambda^2L^2$ and obtains the zero polynomial, with output
\[
\begin{array}{ll}
\texttt{A\_NUMERATOR\_TERMS}&117,\\
\texttt{B\_NUMERATOR\_TERMS}&24,\\
\texttt{CURVE\_NUMERATOR\_TERMS}&248,\\
\texttt{FOURS\_NUMERATOR\_TERMS}&20,\\
\texttt{DIRECT\_NUMERATOR\_TERMS}&0.
\end{array}
\]
No random specialization or floating-point operation enters either identity
check.

The geometric-integrality certificate at $(a,b,c)=(1,2,3)$, modulo $13$,
reports one nonconstant factor of degree eight and exponent one, together
with
\[
 \texttt{POINT\_VALUE}=0,
 \qquad
 \texttt{POINT\_GRADIENT}=(-2,0).
\]
The genuine-open certificate for the second sample reports
\[
\begin{aligned}
\texttt{FIXED\_POINT\_LINE\_VALUES}
 &=(-1,34,19,12,61,40),\\
\texttt{LINE\_SECTION\_DERIVATIVE\_RESULTANTS}&=(-3,6),\\
\texttt{LINE\_BRANCH\_RESULTANTS}&=(-1,1),\\
\texttt{LINE\_PAIR\_RESULTANT}&=-1,
\end{aligned}
\]
 all in $\mathbf F_{13}$ and hence all nonzero.  The source equations,
coefficient formulas, and verification scripts are the supplementary files
\path{code/g2_mu6_source_provenance_verify.py},
\path{code/g2_mu6_lowmem_certificate_fast.py},
\path{code/g2_mu6_direct_numerator_verify.py}, and their inputs in
\path{code/generated/}.  The provenance verifier derives the saved Singular
inputs from the normalized elliptic cubic and the three sections used above,
then compares the resulting formulas and complete files.  The scripts use
exact SymPy arithmetic and can be run with
\begin{verbatim}
python code/g2_mu6_source_provenance_verify.py
python code/g2_mu6_lowmem_certificate_fast.py
python code/g2_mu6_direct_numerator_verify.py
\end{verbatim}
With Singular installed in WSL, the finite-field factorization and
smooth-point checks are reproduced by
\begin{verbatim}
$audit = "code/g2_mu6_universal_special123_mod13_audit.ps1"
powershell -NoProfile -ExecutionPolicy Bypass -File $audit
\end{verbatim}
The complete source and versioned release record are available at
\url{https://github.com/FDmd233/litt-problem-13-g2}.

\bibliographystyle{amsplain}
\bibliography{references}

\end{document}
